\section{Conclusion}
\label{sec:conclusion}

Acoustic perception gives autonomous systems a sense that reaches around corners and through occlusion, and the methods to exploit it, from array processing to large pretrained, audio-visual and language-grounded models, are increasingly mature.
What holds hearing back on real platforms is the platform itself.
This review has argued that ego-noise is the defining problem of the field, and that the platform's own knowledge of its state is the key to solving it.
Across robots and drones, methods that model ego-noise explicitly and use motor or joint data outperform those that treat it as generic interference when ego-noise dominates.
Our case studies on three public datasets quantify this: the advantage is largest exactly when the target is quietest, spatial noise models remove 6 to 13\,dB of ego-noise on all five drones studied, the value of tone-based methods varies widely between drones and flight conditions, and a noise model costs accuracy where there is little structured noise to remove.
Closing the remaining gap will require bringing together three communities that have so far worked largely apart: robot and drone audition, which understands ego-noise; audio machine learning, which provides robust representations; and embodied vision and learning, which provides the tasks, simulators and benchmarks that turn hearing into action.
